\section{Waterfall Cuckoo Details and Alternative Profiles}
\label{app:waterfall-details}

This appendix gives the circuit-level complete-repair generator, the exact
failure calculations behind \tableref{tab:waterfall-parameters}, and the OT
accounting used for its concrete totals.  It then records a shallow
single-repair alternative.  All profiles use the top-level protocol and
simulation theorem from \sectionref{sec:waterfallCuckoo}.

\subsection{Complete-repair circuit}

The circuit represents the current matching by row-to-partition bits
$M_{i,s}$ and the cached candidate unit vectors
$D_{i,s}:=\unitVec_{[0,d_s)}(C_{i,s},1)$.  The physical placement is recovered
as $O_{i,\Delta_s+b}=M_{i,s}D_{i,s}[b]$.  At the start of each repair round,
the circuit materializes from $M$ an occupancy vector $T_s$ and an $h$-bit
owner array $J_s$ for every partition, where $h:=\lceil\log_2t\rceil$.
Thus $T_s[b]=1$ exactly when bin $(s,b)$ is matched, and $J_s[b]$ holds its
matched row.  A secret-index mux reads
$(T_s[C_{i,s}],J_s[C_{i,s}])$.  Every priority choice below is implemented by
the prefix circuit
\[
 \mathsf{first}(x)_i:=x_i\bigwedge_{j<i}(1-x_j).
\]

Let $q_i:=1-\sum_sM_{i,s}$ identify all unmatched rows.  One repair uses a
public $t$-step multi-source reachability search initialized with every row for
which $q_i=1$.  At each step it expands the first visited but unexpanded row in
public row order and reads all $w$ candidate-owner records.  A free record
terminates the search.  An occupied record inserts its owner unless that row
was already visited, and stores the current row and partition as the owner's
parent.  Since a row is expanded at most once, $t$ steps exhaust every row
reachable from any unmatched root.  The circuit then follows the selected
parent chain for $t$ fixed steps and flips the corresponding row-to-partition
assignments.  Gating every write by the selected target makes all unused steps
no-ops.

\begin{figure}[H]\scriptsize
\centering
\fbox{\begin{minipage}{.95\textwidth}
\noindent
$\mathsf{WaterfallGen}_{\mathbf d,c,r}
\bigl(\share C\in\prod_{s\in[0,w)}[0,d_s)^t\bigr)
\longrightarrow\share O\in\{0,1\}^{t\times m}$,
where $m:=\sum_{s\in[0,w)}d_s+c$.
\par\smallskip
Party $p\in\{0,1\}$ computes:

\begin{enumerate}[nosep]
 \item Compute $\share a,\share\ell,\share u$, the initial matching
 $\share M$, and $\share q$ by the occupancy scan in
 \equationref{eq:waterfall-state-transition}.  Cache
 $\share D_{i,s}:=\unitVec_{[0,d_s)}(\share C_{i,s},1)$.

 \item For $v\in[r]$:
 \begin{enumerate}[nosep]
 \item Materialize $(\share T_s,\share J_s)_{s\in[0,w)}$
 from $\share M$ and the cached decoders, and set
 $\share q_i:=1-\sum_s\share M_{i,s}$ for every $i\in[0,t)$.

 \item Initialize $\share V:=\share q$, $\share E:=0^t$, and zero parent
 arrays $\share\pi\in[0,t)^t$ and $\share\eta\in[0,w)^t$.
 Reset the target-selected flag to zero.

 \item For $a\in[0,t)$, let
 $\share f:=\mathsf{first}(\share V\land(1-\share E))$ and add $\share f$
 to $\share E$.  For each $s\in[0,w)$, read the occupancy bit and owner at the
 candidate of $\share f$ from $(\share T_s,\share J_s)$.  Select the first
 free candidate if no target has yet been selected.  Otherwise, insert each
 newly seen owner into $\share V$ in partition order. Store
 $(\sum_{i\in[0,t)}i\share f_i,s)$ in its parent record on first discovery.
 Gate every read and write by an active selected row and an unset target flag.

 \item Starting from the row and partition of the selected free candidate,
 perform $t$ gated backtracking steps.  At each step, assign the current row
 to the current partition, clear its previous partition assignment, and, if
 $\share q_i=0$ for the current row $i$, move to its stored parent record.
 Reaching a root ends the active path; if no free candidate was selected,
 leave $\share M$ unchanged.
 \end{enumerate}

 \item Recompute $\share q_i:=1-\sum_s\share M_{i,s}$ and
 $\share R_i:=\sum_{j<i}\share q_j$.  Set
 $\share O_{i,\Delta_s+b}:=\share M_{i,s}\share D_{i,s}[b]$ and
 $\share O_{i,\Delta_w+k}:=\share q_i[\share R_i=k]$.

 \item Return $\share O$.
\end{enumerate}
\end{minipage}}
\caption{Gate-level organization of the canonical complete-repair generator.
All loops have public bounds; muxes, unit-vector decoders, prefix priority, and
gated writes expand to Boolean products and XORs.
\label{fig:waterfallReachabilityProtocol}}
\end{figure}

\begin{lemma}[Complete-search invariant]
\label{lem:waterfall-reachability-round}
Before the search selects a free bin, $V$ consists exactly of the unmatched
roots and the matched owners discovered from expanded rows.  Every discovered
non-root row stores the first edge by which it was discovered, and the parent
records form a forest rooted at the unmatched rows.  If no free bin has been
selected after $t$ steps, every row reachable from an unmatched root has been
expanded.  If the search selects a free bin, the parent records define a
simple augmenting path from one root to that bin.
\end{lemma}

\begin{proof}
Initially, $V$ contains exactly the unmatched roots.  Expanding a visited row
examines all of its candidate edges.  A free endpoint completes an alternating
path.  An occupied endpoint extends the path through its uniquely matched
owner.  The visited mask inserts each owner once, so its first stored parent is
never overwritten and the parent relation is a forest.  Conversely, the
circuit adds a non-root row only through such an occupied candidate edge.
Induction proves the discovery claim.  Unless a free endpoint terminates the
search, each of the $t$ steps expands one previously unexpanded reachable row;
there are at most $t$ rows.  Backtracking follows one tree of the parent forest
to its unmatched root and therefore flips exactly one simple augmenting path.
\end{proof}

The path flip preserves the row and column inequalities in
\equationref{eq:waterfall-placement-invariants} and increases the main
matching size by one.  Candidate decoders ensure
\equationref{eq:waterfall-placement-subset}.  This proves the invariant
used by \namedref{Proposition}{prop:waterfall-reachability-repair}.

\subsection{Exact failure calculations}

Let $L_s$ be the number of live rows entering partition $s$, with $L_0=t$.
If $L_s=n$ and the $n$ independent candidates occupy $j$ bins, then
$L_{s+1}=n-j$.  Write
$\genfrac\{\}{0pt}{}{n}{j}$ for a Stirling number of the second kind.  Then
\begin{equation}\label{eq:waterfall-occupancy-transition}
 \Pr[L_{s+1}=n-j\mid L_s=n]
 =\frac{(d_s)_j\genfrac\{\}{0pt}{}{n}{j}}{d_s^n}.
\end{equation}
Dynamic programming over this transition gives the exact overflow tail
\begin{equation}\label{eq:waterfall-overflow-tail}
 \Pr[L_w>r+c]=\sum_{\ell>r+c}\Pr[L_w=\ell].
\end{equation}
For a fixed set of $q$ rows, the probability that their candidates occupy
exactly $j$ bins in partition $s$ is
\begin{equation}\label{eq:waterfall-fixed-set-occupancy}
 \Pr[|N_s(Q)|=j]
 =\frac{(d_s)_j\genfrac\{\}{0pt}{}{q}{j}}{d_s^q}.
\end{equation}

Hall's theorem fails only if some $q$ rows have at most $q-1$ total neighbors.
A union bound over the row set and its per-partition neighbor counts gives
\begin{equation}\label{eq:waterfall-hall-bound}
 \Pr[G_C\text{ has no size-$t$ matching}]
 \leq\sum_{q=1}^{t}\binom tq
 \sum_{\substack{j_0,\ldots,j_{w-1}\geq0\\
                  \sum_sj_s\leq q-1}}
 \prod_{s\in[0,w)}
 \frac{(d_s)_{j_s}\genfrac\{\}{0pt}{}{q}{j_s}}{d_s^q}.
\end{equation}
The implementation evaluates \namedeqref{Equations}{eq:waterfall-overflow-tail}
and~\eqref{eq:waterfall-hall-bound} with exact rational arithmetic before
converting the final bounds to bits.

One useful necessary check comes from identical candidate vectors.  Put
$D:=\prod_sd_s$.  If some candidate vector occurs more than $w$ times, those
rows have only $w$ neighbors and violate Hall's condition.  Exact occupancy
counting gives
\begin{equation}\label{eq:waterfall-identical-candidate-probability}
 \Pr[\text{maximum candidate multiplicity}>w]
 =1-\frac{t!}{D^t}[z^t]
 \left(\sum_{j=0}^{w}\frac{z^j}{j!}\right)^D.
\end{equation}
This check rules out overly small candidate spaces before the full Hall sum is
evaluated.

For $t=16$ and $B_0=256$, the complete-repair profiles give
\[
\begin{array}{c|cc}
 & \text{Hall bits} & \text{overflow bits}\\ \hline
(16,16,16,16),\ c=0,r=3 & 43.83 & 45.24\\
(128,128,64),\ c=0,r=2 & 41.15 & 47.07
\end{array}
\]
Combining the two terms in
\equationref{eq:waterfall-reachability-rejection} gives 43.37 and 41.13
batch-correctness bits, respectively.

\subsection{OT accounting}\label{sec:waterfallMpcCost}

We compare profiles by $(Q_{\mathsf{OT}},B_{\mathsf{OT}})$, where the first
coordinate counts OTs and the second sums their payload lengths.  A product of
a shared bit and a shared $\lambda$-bit string costs two $\lambda$-bit OTs and
contributes $(2,2\lambda)$.  XORs, public linear maps, openings, local PRG
calls, depth, and rounds are free.

Let $k_s:=\log_2d_s$, $h:=\lceil\log_2t\rceil$,
$D:=\lceil\log_2N\rceil$, $\ell:=\max\{1,D\}$,
$g:=\log_2|\mathbb G|$, $h_m:=\lceil\log_2m\rceil$, and
$L_{\mathsf{sc}}:=t+D+2$.  To evaluate all hash polynomials, compute powers
of each secret $\iota(A_i)$ once and reuse them across partitions.  Squaring
is linear; only the $\nu_t:=\lfloor(t-2)/2\rfloor$ odd powers require a field
multiplication.

Recursive polynomial-basis Karatsuba gives
\begin{equation}\label{eq:waterfall-karatsuba-cost}
 K(1):=1,
 \qquad
 K(\ell):=K(\lfloor\ell/2\rfloor)+2K(\lceil\ell/2\rceil).
\end{equation}
At $\ell=20$, use
$\mathbb F_{2^{20}}\cong\mathbb F_{16}[Y]/(Y^5+Y^2+1)$ and evaluate the two
degree-four operands at nine points of $\mathbb F_{16}$.  Nine pointwise
$\mathbb F_{16}$ products, each costing $K(4)=9$ bit products, give an
81-product multiplier.  Evaluation, interpolation, reduction, and basis
conversion are binary linear maps.  Put $M(20):=81$ and
$M(\ell):=K(\ell)$ otherwise.  Hashing and dummy selection contribute
\begin{equation}\label{eq:waterfall-hash-ot-cost}
 q_{\mathsf{hash}}:=t\nu_tM(\ell)+wt,
 \qquad
 b_{\mathsf{hash}}:=t\nu_tM(\ell)+t\sum_sk_s.
\end{equation}

For a sparse set $S_j$, let $\sigma(S_j)$ count the binary trie depths that
contain a split, and put
\begin{equation}\label{eq:waterfall-dpf-split-levels}
 \Sigma(P):=\sum_{j:S_j\neq\emptyset}\sigma(S_j)\leq mD.
\end{equation}
Each binary correction selects between shared child sums. It therefore
costs two OTs, not one. We charge $\kappa$ bits per selected message,
including the implementation's full-block representation.
At a sparse level, masking both child sums costs a further two OTs with
$2\kappa$-bit messages. The mask depends on secret activity and cannot be
omitted merely because public splitting nodes exist.

The implementation first expands a complete binary prefix of depth
$D_0:=\min\{D,h_m+2\}$, then prunes the remaining public subtrees.
Activity is known to be one at each complete level; only the residual
sparse levels need masking. The row-to-column trees are fully dense.
For a conservative reservation bound, put
\[
 C_{\mathsf{DPF}}:=mD+th_m,
 \qquad M_{\mathsf{DPF}}:=m(D-D_0).
\]
These count correction selections and activity masks, respectively.
Skipping unused public levels can reduce consumption below this bound.
Let $\chi=0$ for groups of characteristic two and $\chi=1$ otherwise.
In the latter case, each DPF additionally caches one sign--payload product,
costing two OTs with $g$-bit messages. Thus reusable DPF initialization contributes
\begin{align}
 Q_{\mathsf{DPF}}&=2\bigl(C_{\mathsf{DPF}}+M_{\mathsf{DPF}}+\chi(m+t)\bigr),
 \label{eq:waterfall-dpf-ot-count}\\
 B_{\mathsf{DPF}}&=2\kappa C_{\mathsf{DPF}}+4\kappa M_{\mathsf{DPF}}
                  +2g\chi(m+t).
 \label{eq:waterfall-dpf-ot-bits}
\end{align}
The sign controls depend on cached tags, so their correlations are charged
once, not once per expansion. This accounting does not change the leaf
conversion or local expansion work.

If $q_m$ is the number of switches in one recursive
Waksman pass, then
\[
 q_1:=0,
 \quad q_2:=1,
 \quad q_m:=2\lfloor m/2\rfloor+q_{\lfloor m/2\rfloor}
 +q_{\lceil m/2\rceil}\quad(m>2).
\]
The two serial privately controlled passes contribute
\begin{equation}\label{eq:waterfall-waksman-ot-cost}
 (2q_m,2q_mL_{\mathsf{sc}}).
\end{equation}

For deduplication, cache the first-occurrence bits and the equality bits
gated by those first-occurrence bits. Expansion multiplies each retained
payload by its activity bit and routes each later equal payload to its first
occurrence. This requires at most $t+\binom t2=t(t+1)/2$ shared
bit--payload products. Their controls depend only on setup addresses.
The model charges the corresponding OT correlations once, with $g$-bit
payload width; round-separated expansion reuses them but still communicates
masked payloads. The formulas below do not count those online messages.
The common deduplication and activity path has
\begin{align}
 q_{\mathsf{common}}
  &=\binom t2D+\frac{(t-1)(t-2)}2+tD
    +\frac{t(t+1)}2+3t,
    \label{eq:waterfall-common-ot-count}\\
 b_{\mathsf{common}}
  &=\binom t2D+\frac{(t-1)(t-2)}2+t(D+1)
    +\frac{t(t+1)}2g+t+t(2D+1).
    \label{eq:waterfall-common-ot-bits}
\end{align}

The optimized initial scan performs one mux-tree occupancy read, one gated
unit-vector decoder, and one live-bit product per row and partition.  Stash
compaction adds one product per row.  Thus the initial stage has
\begin{align}
 q_{\mathsf{basic}}&=t\sum_s(2k_s+1)+t,
 \label{eq:waterfall-basic-ot-count}\\
 b_{\mathsf{basic}}&=t\sum_s(2d_s-1)+t(c+1).
 \label{eq:waterfall-basic-ot-bits}
\end{align}
For complete repair, the scan caches raw and winning decoders and gates the
final matching once.  Its base profile is
\begin{equation}\label{eq:waterfall-bounded-base-cost}
 q_{\mathsf{base}}=t\sum_s(2k_s+2)+t,
 \qquad
 b_{\mathsf{base}}=t\sum_s(4d_s-3)+t(c+1).
\end{equation}

With packed control bits, one complete augmentation contributes
\begin{align}
 q_{\mathsf{reach}}
  :={}&wt+t\sum_sk_s+t(t+2h+wh+3w+2)\nonumber\\
     &+(w-1)+(t-1)(2h+w)+t,
 \label{eq:waterfall-reachability-augmentation-ot-count}\\
 b_{\mathsf{fwd}}
  :={}&t\bigl((t-1)(2+w(h+4))+ht+4w\bigr),\nonumber\\
 b_{\mathsf{back}}
  :={}&(w-1)(t-1)
     +(t-1)\bigl((h+2w+1)(t-1)+1\bigr)+tw,\nonumber\\
 b_{\mathsf{reach}}
  :={}&t\sum_s(d_s-1)+t(h+1)\sum_s(d_s-1)
     +b_{\mathsf{fwd}}+b_{\mathsf{back}}.
 \label{eq:waterfall-reachability-augmentation-cost}
\end{align}
Consequently,
\begin{align}
 q_{\mathsf{BA}}&=q_{\mathsf{base}}+rq_{\mathsf{reach}},
 \label{eq:waterfall-reachability-ot-count}\\
 b_{\mathsf{BA}}&=b_{\mathsf{base}}+rb_{\mathsf{reach}}.
 \label{eq:waterfall-reachability-ot-bits}
\end{align}
Combining all components gives
\begin{align}
 Q_{\mathsf{OT}}
  &=2(q_{\mathsf{common}}+q_{\mathsf{BA}}+q_{\mathsf{hash}})
    +2q_m+Q_{\mathsf{DPF}},
    \label{eq:waterfall-total-ot-count}\\
 B_{\mathsf{OT}}
  &=2(b_{\mathsf{common}}+b_{\mathsf{BA}}+b_{\mathsf{hash}})
    +2q_mL_{\mathsf{sc}}+B_{\mathsf{DPF}}.
    \label{eq:waterfall-total-ot-bits}
\end{align}
At $t=16$, $N=2^{20}$, $\kappa=128$, and $\mathbb G=\mathbb Z_{2^{64}}$,
these equations give
$(39{,}052,942{,}888)$ for the default $4N$ profile and
$(55{,}358,3{,}623{,}410)$ for the lower-expansion $3N$ profile, matching
\tableref{tab:waterfall-parameters}.
Their dense-prefix depths are 8 and 11. The DPF terms are, respectively,
$(4{,}448,755{,}712)$ and $(19{,}520,3{,}192{,}832)$.
For a characteristic-two group of the same size, omit $2(m+t)$ OTs and
$128(m+t)$ payload bits from each total.

\subsection{Specialized single-repair profile}
\label{app:waterfallOneRepair}

This setup-minimized alternative performs one fixed depth-two augmentation
after the initial waterfall stage.  If $L:=\sum_iq_i\leq c$, the stash holds
all overflow rows.  If $L=c+1$, it selects the first overflow row, inspects the
later candidates of the rows occupying its candidates, and applies the first
free length-two path.  Otherwise the excess rows remain unplaced.  The
protocol opens no failure bit; the simulation theorem charges this event to
$p_{\mathsf{1rep}}$.

\begin{figure}[H]\scriptsize
\centering
\fbox{\begin{minipage}{.95\textwidth}
 \noindent
 $\textsf{OneRepairWaterfallGen}_{w,d,c}
 \bigl(\share C\in[0,d)^{t\times w}\bigr)
 \longrightarrow\share O\in\{0,1\}^{t\times(wd+c)}$.
 \par\smallskip
 Party $p\in\{0,1\}$ computes:
\begin{enumerate}
  \item Compute $\share a,\share\ell,\share u$, the main-position matrix
  $\share O^{(0)}$, $\share q$, and $\share R$ as in
  \namedeqref{Equations}{eq:waterfall-state-transition},
  \eqref{eq:waterfall-main-owner}, and~\eqref{eq:waterfall-stash-rank}.

  \item $\share L:=\sum_{i\in[0,t)}\share q_i$,
  $\share x_i:=\share q_i\land\textsf{eq}(\share R_i,0)$ for $i\in[0,t)$,
  and $\share f_{\mathsf{try}}:=\textsf{eq}(\share L,c+1)$.

  \item For $u\in[0,w)$, let
  $\share b_u:=\sum_{i\in[0,t)}\share x_i\share C_{i,u}$ and, for $i\in[0,t)$,
  \[
    \share v_{u,i}:=\bigvee_{b\in[0,d)}
      \bigl(\textsf{eq}(\share b_u,b)\land\share O^{(0)}_{i,ud+b}\bigr).
  \]
  Write $\share v_u$ for the row selected by the one-hot vector
  $(\share v_{u,i})_{i\in[0,t)}$, and similarly write $\share x$ for
  $(\share x_i)_{i\in[0,t)}$.

  \item For $(u,s)\in[0,w)^2$ with $u<s$, let
  $\share b_{u,s}:=\sum_{i\in[0,t)}\share v_{u,i}\share C_{i,s}$ and
  \[
    \share g_{u,s}:=\share f_{\mathsf{try}}\land
      \bigwedge_{i\in[0,t)}\bigwedge_{b\in[0,d)}
      \neg\bigl(\textsf{eq}(\share b_{u,s},b)
                 \land\share O^{(0)}_{i,sd+b}\bigr).
  \]

  \item In lexicographic order on $(u,s)$, let
  $\share z_{u,s}:=\share g_{u,s}\land
  \neg\bigvee_{(u',s')<(u,s)}\share g_{u',s'}$ and
  $\share z:=\bigvee_{u<s}\share z_{u,s}$.

  \item For $(i,a,b)\in[0,t)\times[0,w)\times[0,d)$, set
  \begin{align*}
   \share\Delta^{\mathsf{old}}_{i,a,b}
    &:=\bigoplus_{u<s}\share z_{u,s}[a=u]
       \textsf{eq}(\share b_u,b)
       (\share v_{u,i}\oplus\share x_i),\\
   \share\Delta^{\mathsf{new}}_{i,a,b}
    &:=\bigoplus_{u<s}\share z_{u,s}[a=s]
       \textsf{eq}(\share b_{u,s},b)\share v_{u,i},\\
   \share O^{(1)}_{i,ad+b}
    &:=\share O^{(0)}_{i,ad+b}
       \oplus\share\Delta^{\mathsf{old}}_{i,a,b}
       \oplus\share\Delta^{\mathsf{new}}_{i,a,b}.
  \end{align*}

  \item For $i\in[0,t)$, let
  $\share q'_i:=\share q_i\land\neg(\share x_i\land\share z)$ and
  $\share R'_i:=\sum_{j<i}\share q'_j$.  For $(k,i)\in[0,c)\times[0,t)$, set
  $\share O^{(1)}_{i,wd+k}:=\share q'_i\land
  \textsf{eq}(\share R'_i,k)$.

  \item Return $\share O^{(1)}$.
\end{enumerate}
\end{minipage}}
\caption{The specialized single depth-two repair followed by stash
compaction.\label{fig:waterfallOneRepairProtocol}}
\end{figure}

\subsection{Failure bound and parameters}

For a live-count path $\boldsymbol L=(L_0,\ldots,L_w)$ under
\equationref{eq:waterfall-occupancy-transition}, put
$J_s:=L_s-L_{s+1}$ and let $p(\boldsymbol L)$ denote the path probability.

\begin{proposition}[One-repair failure bound]
\label{prop:waterfall-one-repair}
Fix $d\in\mathbb N$, and define
\[
 \mathbf d:=(d,\ldots,d)\in\mathbb N^w.
\]
Then
\begin{equation}\label{eq:waterfall-one-repair-bound}
 p_{\mathsf{1rep}}
 \leq \Pr[L_w\geq c+2]
 +\sum_{\boldsymbol L:L_w=c+1}
   p(\boldsymbol L)
   \prod_{s=1}^{w-1}\left(\frac{J_s}{d}\right)^s.
\end{equation}
\end{proposition}

\begin{proof}
The first term covers every proposal with at least $c+2$ overflow rows.  Fix a
path ending in $c+1$ rows and let $x$ be the selected
overflow row.  Its candidate in partition $u$ is occupied by a row $v_u$
placed in that partition.  The $v_u$ are distinct because a placed row owns
only one position.  Every candidate of $v_u$ through partition $u$ was tested
by the waterfall and is occupied.  For every later partition $s>u$, however,
$C_{v_u,s}$ was never inspected because $v_u$ had already left the live set.
These cells are mutually independent and uniform, and are independent of the
$J_s$ occupied bins in partition $s$.  Exactly $s$ such untouched cells target
partition $s$, one from each $v_0,\ldots,v_{s-1}$.  All length-two paths are
blocked with probability $\prod_{s=1}^{w-1}(J_s/d)^s$.  Summing over the
live-count paths proves the bound.
\end{proof}

For $t=16$ and $B_0=256$, the compact proved profile is:

\begin{table}[H]
\centering
\small
\begin{tabular}{@{}rrrrrrrr@{}}
\toprule
$w$ & $d$ & $c$ & $m$ & expansion & total OTs & OT payload & batch bits \\
\midrule
4 & 16 & 2 & 66 & $6N$ & 30{,}980 & 829{,}288 bits & 45.23 \\
\bottomrule
\end{tabular}
\caption{The retained specialized one-repair profile.
\label{tab:waterfall-one-repair-parameters}}
\end{table}

The compact point $(w,d,c)=(4,16,2)$ satisfies
\begin{align}
 \Pr[L_4\geq4]&<2^{-53.2360},\nonumber\\
 \Pr[L_4=3\text{ and depth-two repair fails}]&<2^{-60.6651},\nonumber\\
 p_{\mathsf{1rep}}&<2^{-53.2277}.
 \label{eq:waterfall-compact-point}
\end{align}
Thus $B_0p_{\mathsf{1rep}}<2^{-45.2277}$.  It uses 66 position columns and $6N$
local expansion.

\subsection{OT count}

Put $k:=\log_2d$, $h:=\lceil\log_2t\rceil$,
$\omega:=\lceil\log_2w\rceil$, and $P:=\binom w2$.  For the OT-minimized
realization, augment the occupancy scan from
\sectionref{sec:waterfallMpcCost} with an $h$-bit owner array
$J_s\in[0,t)^d$ for each partition.  Whenever row $i$ wins with gated unit vector
$E_i$, update the bit planes of $J_s$ by the public binary encoding of $i$
times $E_i$.  This is linear.  Reading $J_s[C_{i,s}]$ later returns the owner
of an occupied candidate without an all-pairs equality scan.

During the scan, also cache the ungated unit vector for every $C_{i,s}$.
Its decoder shares the same address bits and dependency level as the winning
unit-vector decoder, so their string operands are concatenated in each OT.
After the repair, one wide product gates the cached vector by the final
row-to-partition matching bit; one coordinate is recovered linearly from the
unit-vector sum.  Together, the scan, both decoders, and final gate use
$2k+2$ products with width sum $4d-3$ per row and partition.

The one-hot overflow root gives its $h$-bit row number by a public linear map.
One mux tree, batched across all $w$ candidate arrays, reads its candidates.
The $w$ owner reads use the corresponding $k$-bit candidate addresses.  For a
fixed victim, reads of all its later-partition candidates share the same row
address and are likewise batched.  The resulting $P$ candidate addresses read
the occupancy arrays.  Finally, a priority scan selects the first feasible
pair, a mux selects its victim identifier, and two demux trees route the
$w$-bit matching deltas to the root and victim rows.  This gives
\begin{align}
 q_{\mathsf{1rep}}
  ={}&wt(2k+2)+2t+(w+2)h+k(w+P)+P+\omega+3,
  \label{eq:waterfall-one-repair-ot-count}\\
 b_{\mathsf{1rep}}
  ={}&wt(4d-3)+t(2c+3)+wk(t-1)+wh(d-1)\nonumber\\
    &+Pk(t-1)+P(d-1)+2P+2^\omega h+2w(t-1)+1.
  \label{eq:waterfall-one-repair-ot-bits}
\end{align}
At $(t,w,d,c)=(16,4,16,2)$ these are
$q_{\mathsf{1rep}}=747$ and $b_{\mathsf{1rep}}=5{,}095$.
These expressions are conservative upper bounds; omitting the unused explicit
failure test only lowers them.  Including the common, direct-hash, scatter,
column sparse-DPF, and cached row-to-column DPF terms from
\namedeqref{Equations}{eq:waterfall-total-ot-count}
and~\eqref{eq:waterfall-total-ot-bits} gives at most $30{,}980$ OTs carrying
$829{,}288$ payload bits for $\mathbb G=\mathbb Z_{2^{64}}$.

\subsection{Asymmetric complete-repair profile}

The complete-repair generator also supports
$\mathbf d=(32,32,16,8)$, $c=0$, and $r=2$.  This profile has $m=88$,
$41.20$ bits of batch correctness, and $4N$ local expansion.  Substitution in
\namedeqref{Equations}{eq:waterfall-total-ot-count} and
\eqref{eq:waterfall-total-ot-bits} gives $38{,}158$ OTs carrying $1{,}162{,}710$
payload bits. It saves $894$ OTs relative to the default profile, but adds
24 physical columns, about $26.8$ KiB of OT payload, and a smaller correctness
margin.  We therefore retain it as a setup-versus-state alternative rather
than a main profile.

\subsection{Hybrid simulator for Waterfall}
\label{app:waterfall-simulator}

This subsection supplies the simulator used in
\namedref{Theorem}{thm:waterfall-simulation}.  Fix a static semi-honest corrupted
party $p\in\{0,1\}$.  Circuit, permutation, and reusable DPF calls are ideal.
The simulator receives party $p$'s setup input shares and, for each expansion,
its payload shares. It maintains the simulated caller's shares and public
state between calls. Ideal subprotocols expose no additional implementation
state. At every ideal DPF expansion, the corrupt party selects its output
share using its current view. The simulator runs this strategy on the
simulated view; it does not replace the selected map by an independent sample.
When composing concrete DPF realizations, their stateful simulators supply
these strategies and simulate their own local states and messages.

\paragraph{Setup calls and opening.}
The simulator samples each public $h_s$ from $\mathcal H_{t,k_s}$ and derives
$P$ and $S$.  It follows the setup circuit in order.  For each fresh secret
output of an ideal circuit call, it samples the corrupted party's share with
the prescribed uniform law.  It computes local linear operations, public
constants, and repeated uses of a share rather than resampling them.
This simulates the normalization and placement calls without knowing the
honest party's address shares.

Each ideal permutation pass returns freshly shared outputs while retaining
its private permutation for the inverse call.  The simulator samples the
corrupted party's output shares for these calls. It also retains the corrupt
party's sampled permutation and its inverse, which are part of that party's
local view. At the opening, it samples
$V\gets\mathsf{Inj}(t,m)$ and sends the honest opening share
$V\oplus\widehat O_p$.  The public scatter is then a deterministic operation
on $V$ and the previously sampled shares.  In particular, the simulator
does not independently resample the public scatter's outputs.  It continues
with the inverse permutation outputs, computes the representative addresses,
and supplies the setup responses of the row and column DPF functionalities.
On successful routing their indices are valid, so those responses are
$\mathsf{ready}$.

For every fixed hash tuple and successful candidate matrix, the hidden
honest permutation makes the real opening uniform on $\mathsf{Inj}(t,m)$.
Couple that opening to $V$.  The simulated hashes are still unconditioned
samples: routing success need not be independent of the hashes.  If routing
fails, discard this coupling and charge the entire continuation to the
failure event.  This includes the unmatched row labels exposed by the opening.
Its probability is $p_\Gamma$ for every fixed setup input.

\paragraph{Expansion calls.}
The simulator retains the setup shares used for duplicate aggregation.
It samples the fresh private outputs of the ideal multiplication calls and
computes their specified local sums. For each row $i\in[0,t)$, let
$v_{p,i}\in\mathbb G^{[0,m)}$ be the map chosen by the corrupt party at its
row-to-column DPF call. The simulator computes the corrupt column-payload
share $B'_{p,j}:=\sum_{i\in[0,t)}v_{p,i}(j)$ for each $j\in[0,m)$.
These are the inputs to its column-DPF calls. No row destination, equality
bit, or setup state is sampled again.

For each nonempty $S_j$, let $Z_{p,j}\in\mathbb G^{S_j}$ be the corrupt
output map selected at column DPF $j$. The simulator follows those selections
and computes
\begin{equation}\label{eq:waterfall-sim-output-completion}
 Y_p(x):=\sum_{j:x\in S_j}Z_{p,j}(x),\qquad x\in[0,N).
\end{equation}
It submits $Y_p$ to $\mathcal F_{\mathsf{rDMPF}}$. Empty sparse sets require
no DPF call. This procedure imposes no distribution on the selected maps
beyond that of the corrupt strategy on its simulated view.

\paragraph{Honest outputs and reuse.}
On successful routing, each ideal DPF gives the honest party the complement
to the specified point map. Summing those complements and applying
\namedref{Proposition}{prop:waterfall-correctness} gives the honest DMPF output
$\mathsf{mp}_{N,t}(A,B)-Y_p$. This is exactly the output of the top-level
functionality. Thus the coupling preserves the joint corrupt view and honest
output, not merely the marginal transcript. It holds for arbitrary choices
of the corrupt subprotocol shares, including correlated, nonuniform choices.

Given a coupled history, use the same next inputs and corrupt random coins
in both executions. Fresh circuit-output shares can be coupled identically;
DPF share selections then follow the same strategy on the same view.
The complement identity completes the honest output for the next call.
Induction proves the hybrid claim for polynomially many adaptive expansions.
Routing failure is fixed during setup and is charged once per state.
Replacing concrete subprotocol calls by their ideal calls and online
simulators adds their accumulated joint simulation loss
$\epsilon_{\mathsf{sub}}(B_0,Q)$. This loss includes any distinguishable
failure of those realizations; the hybrid routing error remains separate.
